% ----------------------------------------------------------------------
% Self-Given Is Not Self-Grounding
% Copyright (c) 2026 Zain Dana Harper. Written by Zain Dana Harper.
% License: CC BY 4.0, https://creativecommons.org/licenses/by/4.0/
% DOI: 10.5281/zenodo.23126458
% First public: 2026-09-15 (Zenodo, first version (10.5281/zenodo.22768919))
% Source SHA-256: 50cc0e469edb09a4538ed31f7f06ee6e30d57fa9dfef08e4b94a385f862dff39 (writing/papers/self-given.md)
% ----------------------------------------------------------------------
\documentclass[11pt]{article}
\usepackage[a4paper,margin=1in]{geometry}
\usepackage{fontspec}
\setmainfont{lmroman10-regular.otf}[BoldFont=lmroman10-bold.otf,ItalicFont=lmroman10-italic.otf,BoldItalicFont=lmroman10-bolditalic.otf,Ligatures=TeX]
\setmonofont{lmmono10-regular.otf}
\usepackage{microtype}
\usepackage{longtable,array,booktabs}
\usepackage{enumitem}
\usepackage[hidelinks]{hyperref}
\usepackage{url}
\setlist{itemsep=2pt,topsep=4pt}
\setlength{\parskip}{6pt plus 1pt}
\setlength{\parindent}{0pt}
\setlength{\emergencystretch}{3em}
\hypersetup{pdftitle={Self-Given Is Not Self-Grounding},pdfauthor={Zain Dana Harper},pdfsubject={Typeset from writing/papers/self-given.md},pdfkeywords={source sha256 50cc0e469edb09a4538ed31f7f06ee6e30d57fa9dfef08e4b94a385f862dff39}}
\newcommand{\register}[1]{{\ttfamily\small #1}}
\usepackage{xcolor}
\makeatletter
\def\ps@attribution{\let\@mkboth\@gobbletwo\let\@oddhead\@empty\let\@evenhead\@empty\def\@oddfoot{\parbox[b]{\dimexpr\textwidth-3em\relax}{\raggedright\scriptsize\color{black!60}\textcopyright{} 2026 Zain Dana Harper \textperiodcentered{} CC BY 4.0 \textperiodcentered{} doi.org/\allowbreak{}10.5281/\allowbreak{}zenodo.23126458 \textperiodcentered{} first public 15 September 2026}\hfill{\small\thepage}}\let\@evenfoot\@oddfoot}
\let\ps@plain\ps@attribution
\makeatother
\pagestyle{attribution}
\hypersetup{pdfauthor={Zain Dana Harper}}
\AtBeginDocument{\special{pdf:docinfo<</Copyright (Copyright 2026 Zain Dana Harper. Licensed under CC BY 4.0, https://creativecommons.org/licenses/by/4.0/)>>}}
\begin{document}

\begin{center}
{\LARGE\bfseries Self-Given Is Not Self-Grounding\par}
\medskip
{\large\emph{Perspectival aseity and its strongest contraries}\par}
\medskip
Zain Dana Harper\par
{\small Philosophy paper · edition of 3 October 2026\par}
{\small Argument built June 2026; first published 15 September 2026; this edition 3 October 2026.\par}
\end{center}

\begin{quote}\small
\textbf{Revised 3 October 2026.} This edition replaces the edition first published on 15 September 2026, in its 25 September 2026 web revision. The paper is rebuilt as a series of contests in which each position meets its strongest contrary, and its passages are marked plain, steelman or open. The earlier claim that the everyday subject dissolves into a dependent occurrence is narrowed: no argument examined establishes its aseity, and whether causal conditioning makes it dependent for its existence is left open. The lamp argument is now credited to Nāgārjuna, and Candrakīrti's target is corrected. Descartes, Strawson and Zahavi are recorded as allies on aseity, and Miri Albahari's view is added, so seity no longer includes ownerlessness. Seven open questions are handed to the reader. A dated update, October 2026, engages Julian Jaynes. This edition is deposited at Zenodo, \url{https://doi.org/10.5281/zenodo.23126458}. The prior version, the edition of 15 September 2026, stays citable at its Zenodo record, \url{https://doi.org/10.5281/zenodo.22768919}.
\end{quote}
\section*{How to read this paper}

The paper sets one idea against the strongest positions that oppose it, and reports how each contest ends. Every section carries one of three marks.

\begin{itemize}
\item \register{[PLAIN]} marks what can be said flatly: a record that was checked, a derivation that is valid, a defeat the paper takes, or a null result. These passages carry no hedge.
\item \register{[STEELMAN]} marks a position set against its strongest contrary. Both sides are argued at a strength their holders should recognise. Each contest ends with one of four outcomes: \emph{defeat}, \emph{survival by concession}, \emph{standoff}, or \emph{ally mistaken for an enemy}.
\item \register{[OPEN]} marks a question handed to you. It is framed as a question, with the considerations on each side and what would settle it.
\end{itemize}
Every numbered claim also carries a domain tag. \emph{C-M} marks a metaphysical claim about consciousness (what exists, what depends on what). \emph{C-E} marks an empirical claim about consciousness. \emph{X} marks a report of what a text or tradition holds, or a rule of method; the claim it reports keeps its own tag. \emph{T} (theological) and \emph{E} (ethical) claims are named where they appear, and no inference crosses from one domain to another without a bridge stated in its own sentence.

\section*{Abstract}

\register{[PLAIN]} The paper grants, as a premise, that experience is known from itself; section 10.2 records who denies it. Does it follow that experience exists from itself? This paper tests that step in five places where it has been argued or attributed: Descartes's cogito, Galen Strawson's real phenomenal subject, the doctrine of self-luminosity (\emph{svaprakāśa}) as developed through Dan Zahavi, Fichte's self-positing I, and the self-established witness of Advaita Vedānta. A sixth contrary, Miri Albahari's ownerless and unconditioned consciousness, is added in this edition. The distinction that decides each case is between self-givenness, an epistemic property (a thing known or warranted from itself), and self-grounding or aseity, an ontological property (an existence that is ungenerated and depends on nothing else).

The results differ by case, and the paper reports each one as it fell. Four arguments establish self-givenness and stop short of aseity. Three of their authors (Descartes, Strawson, Zahavi) decline the step to aseity themselves, so on that question they are allies of the thesis more than opponents. Fichte's absolute I makes the step and does not earn it. Advaita makes the step at the ultimate level and stands undefeated there, in a standoff with Madhyamaka that this paper does not settle. Albahari's view, which this edition reports and does not examine, shows that an ownerless consciousness need not be held dependent, which corrects how the original edition defined its own conclusion.

What survives plainly is narrower than the original edition claimed. No argument examined here establishes aseity for the everyday first-person subject. Whether that subject is positively dependent for its existence rests on a further premise, and whether the ultimate ground of experience is a self-luminous aseity is left open. The paper says what would settle each open question.

\emph{Keywords:} aseity, seity, self-givenness, self-luminosity, cogito, Advaita Vedānta, Madhyamaka, ownerless consciousness

\section*{1. The question}

\register{[OPEN]} Is consciousness the one thing that exists from itself?

The temptation to say yes has reasons behind it. First-person experience presents itself with unusual force. From the inside it cannot be doubted the way external objects can. It seems to carry its own warrant, and to be in some sense its own evidence. The claim runs from Descartes's \emph{cogito ergo sum} through Fichte's absolute self-positing, and from the Vedāntic witness (\emph{sākṣin}) to the pre-reflective self-awareness that phenomenologists study. Each time it takes a recognisable shape: the conscious subject, or consciousness itself, is held to be the one entity whose being hangs on nothing else.

This paper calls that claim \emph{perspectival aseity}. The question arose in June 2026, when the author asked whether the thesis that nothing has aseity should be given up at its end if perspectival aseity turned out to be true, and whether the truth on the other side, if it is there, should be shown. The paper takes that question at its word. A thesis that cannot lose is not being tested, so each contrary below is argued as its holders would argue it, and the thesis is allowed to fail where it fails.

You are invited to keep score. The paper reports its own verdicts and labels them. Where a verdict depends on a premise you may reject, the paper says which premise it is.

\section*{2. Terms, stated once}

\register{[PLAIN]} These definitions fix how the paper uses its words. They settle nothing by themselves.

\textbf{Self-givenness} is an epistemic property. A thing is self-given when it is known or warranted from itself: it needs no outside certification to count as manifest, its presence cannot be doubted in some characteristic way, and its character is intrinsic: relations to other things make up no part of it. Self-luminosity in the Indian traditions, Sydney Shoemaker's immunity to error through misidentification (Shoemaker 1968) and Zahavi's pre-reflective for-me-ness (Zahavi 2005) are all forms of it. The paper grants these phenomena as premises and does not argue for them. Illusionists deny phenomenal for-me-ness as ordinarily conceived (Frankish 2016), and Marie Guillot argues that for-me-ness, me-ness and mineness are distinct and often run together (Guillot 2017). Section 10.2 takes up what follows for the paper if they are right.

\textbf{Aseity}, from the Latin \emph{a se}, ``from oneself'', is an ontological property. A thing has aseity when it exists from and through itself: it is ungenerated, and its existence is not conferred by, derived from or sustained by anything else. Classical theism gives aseity to God alone; Spinoza gives it to the one substance. The paper also says \textbf{self-grounding} for this property. It uses that word as a synonym for aseity only. In current metaphysics grounding is usually treated as irreflexive (Schaffer 2009), and read in that idiom ``self-grounding'' would be incoherent by definition; the paper makes no claim about the grounding relation.

\textbf{Dependence} has three senses in what follows, and the paper keeps them apart.

\begin{itemize}
\item \emph{Causal conditioning of occurrence}: a thing occurs only when certain conditions hold, as a flame burns only while there is fuel.
\item \emph{Generation}: a thing begins to exist, and so was once not there.
\item \emph{Existential dependence}: a thing's existence is from another, so that it could not exist unless that other did. Current metaphysics treats this as one of several distinct relations of ontological dependence, separate from causal dependence (Tahko and Lowe, ``Ontological Dependence'', \emph{Stanford Encyclopedia of Philosophy}).
\end{itemize}
Aseity excludes generation by definition, since an a se being is ungenerated. Aseity also excludes existential dependence by definition. Whether causal conditioning of occurrence amounts to existential dependence is a substantive question, taken up in section 6.

\textbf{Seity} names what is left when aseity is denied to something that is still not determined from outside in the plain way a stone is: it is self-given, it occurs, and it is not a se. The original edition also built ``ownerless'' into seity. Section 5.6 shows why this edition separates the two.

\textbf{The no-aseity thesis} comes in two forms. Its \emph{bracketed} form says that nothing in the conferred or conventional order has aseity. Its \emph{universal} form extends that to the ultimate level, which in this paper means the ultimate ground of experiencing. Neither form is a claim about God (BR-1). The \textbf{conferred or conventional order} is the order of things that arise in experience and in the world the sciences describe. That definition builds no dependence in, so the bracketed thesis is a claim that can fail.

\textbf{Two bridges.} The paper makes no theological claim and no ethical one, and two stated premises keep it that way. A reader may reject either on its own.

\begin{itemize}
\item \textbf{BR-1.} ``Aseity'' names a metaphysical property, existence from nothing else, which can be stated and tested with no claim about God. Classical theism and Spinoza appear here only as reports of where the property has been placed. Nothing in the paper bears for or against theism, or on whether God has aseity.
\item \textbf{BR-2.} Claims about Brahman, Ātman, \emph{turīya} and \emph{nirvāṇa} are reported as the metaphysical content of philosophical arguments. The paper takes no position on their religious or soteriological truth.
\end{itemize}
A third bridge, BR-3, is needed where an empirical premise is used for a metaphysical conclusion. It is stated in section 6, where it does its work.

\section*{3. The flame, and what it can show}

\subsection*{3.1 The image}

\register{[PLAIN]} Intrinsic character does not entail aseity of existence. One counterexample is enough to establish that, and a clean one is a statue. Its shape belongs to it and is no relation to anything outside it, and the statue exists only while its clay does. \emph{C-M.}

The original edition made the point with a flame, and the paper keeps the flame as its picture. Heat seems part of what being a flame is, yet the flame needs fuel; take the fuel away and the flame loses its existence, which is more than losing a property. As an illustration this works. As a counterexample it can be contested, since a physicist may say a flame's heat is released by its reaction with the fuel, and ``heat'' in physics names energy passing between bodies. The entailment point rests on the logic and on the statue, and the flame only pictures it. \emph{Illustration from ordinary description.}

\subsection*{3.2 What the image cannot do}

\register{[PLAIN]} The flame tests one inference: from intrinsic character to aseity. It does not test the other two inferences the traditions use. The cogito argues from indubitability, and Advaita argues from ineliminability and non-objectifiability. A flame's heat is known by no one from itself, so calling the heat ``self-given'', as the original edition did, works by analogy only. And a fact about combustion cannot show that consciousness depends on anything. The image shows that two properties can come apart. Each case below gets the test that fits its own inference.

\subsection*{3.3 The image against its classical source}

\register{[STEELMAN]} \emph{For the image.} The fire and fuel picture is the plainest model of a thing whose character is its own and whose existence is borrowed. It needs no theory of consciousness to work.

\emph{Against it.} The image has a classical treatment the original edition never cited. Nāgārjuna examines fire and fuel in chapter 10 of the \emph{Mūlamadhyamakakārikā}, the examination of fire and fuel (\emph{agnīndhanaparīkṣā}), and applies the analysis to the self and the aggregates it appropriates. On the summaries checked for this edition, the chapter argues that fire is neither identical with its fuel nor other than it, and that identity and difference here hold only conventionally. [Chapter title and its application to the self checked 2026-10-02 at search level; the verse text was not opened.] A Mādhyamika will press this. The plain reading ``fire depends on fuel'' treats fire and fuel as two things with their own natures, one leaning on the other, and that is exactly the picture chapter 10 takes apart.

\emph{Outcome: survival by concession.} The image survives for the one job section 3.1 gives it, because that job needs only the conventional reading: at the level of ordinary description, a flame's heat is intrinsic and its existence requires fuel. The image cannot be used as a picture of how things finally are. A reader who holds the Madhyamaka view will find the paper's metaphysics shallower than chapter 10's, and on that point the paper concedes the depth.

\section*{4. The argument in steps}

\register{[PLAIN]} The steps below are the original edition's, with the changes of this edition marked. Labels in italics are the paper's own status words, or ``no label'' where it gives none.

\begin{itemize}
\item \textbf{P1.} Self-givenness and self-grounding are distinct properties. Epistemic independence is compatible with ontological dependence. \emph{No label. C-M.} (3.1)
\item \textbf{P2.} Intrinsic character does not entail aseity of existence. \emph{No label. C-M.} (3.1)
\item \textbf{P3.} In each of the five traditions, an argument is available that moves from self-givenness toward self-grounding. \emph{No label. X.} (Section 5. Changed: the original said each tradition makes the move ``in some form''; three of the five authors decline it.)
\item \textbf{P4.} Aseity is ungenerated. A thing that begins to exist lacks aseity. \emph{No label. C-M; true by the definition in section 2.}
\item \textbf{P4a.} The everyday first-person subject, taken as the embodied person of one life, began. \emph{No label. C-E.} (Section 6. Added in this edition: the generation route uses it, and Advaita, which holds the individuated self beginningless, rejects it for that self.)
\item \textbf{P5.} The occurrence of the everyday first-person subject's experience is causally conditioned by neurological processes, developmental history and ongoing causal relations. \emph{No label. C-E.} (Section 6)
\item \textbf{P6.} Śaṅkara's step from the witness's ineliminability and non-derivability to its self-establishment, as reconstructed here, uses a grounding excluded middle: everything is either other-grounded or self-grounded. \emph{No label. X, reporting C-M.} (5.5)
\item \textbf{P7.} Madhyamaka argues for a third option, groundless dependent arising (\emph{śūnyatā}), and contemporary scholarship defends its coherence (for example, Westerhoff 2009). \emph{No label. X, reporting C-M.} (5.5)
\end{itemize}
From these:

\begin{itemize}
\item \textbf{I1.} In the cases of 5.1 to 5.4, what is established is self-givenness, and no argument examined earns the step to aseity. \emph{C-M.}
\item \textbf{I2.} If P7's third option is coherent, Śaṅkara's excluded middle begs the question, and the ultimate level is a standoff. \emph{C-M.}
\item \textbf{C1, defensive form (this edition).} No argument examined here establishes aseity for the everyday first-person subject. \emph{C-M. Stated flat; see section 7.}
\item \textbf{C1, dependence form.} The everyday first-person subject lacks aseity because it is generated (P4a with P4), and its experience is causally conditioned (P5). The generation route is valid if the subject begins to exist; the conditioning route needs BR-3. \emph{C-M from C-E: the generation route through P4, which works as its bridge; the conditioning route through BR-3. Status: section 6.}
\item \textbf{C2.} Whether the ultimate ground of experiencing is a self-luminous aseity is open, bracketed and unsettled. The no-aseity thesis cannot be maintained in its universal form. \emph{Open question, bracketed and unsettled. C-M.}
\end{itemize}
The original edition labelled its empirical conclusion ``decisive (high confidence)''. Section 7 explains why this edition states the conclusion in two parts.

\section*{5. Six contests}

\subsection*{5.1 The cogito against Lichtenberg and Kant}

\register{[STEELMAN]} \emph{For the cogito.} You cannot doubt that you are doubting without doing some doubting. Doubting is thinking, and wherever there is thinking, the argument runs, there is an I that thinks. So \emph{I am} cannot be overturned. Descartes also gives a reason why the thinking must have an owner: thought is an attribute, nothing has no attributes, so wherever an attribute is perceived there is a substance that bears it (\emph{Principles} I.11 and I.52; UNVERIFIED at the locus). The I appears self-given, and it seems to secure its own existence through the very attempt to doubt it away. Its strongest modern reconstructions add force. Shoemaker shows that first-person reports cannot go wrong about whose experience is being reported (Shoemaker 1968). Thomas Nagel shows that experience has a subjective character, something it is like, which third-person description leaves out (Nagel 1974).

\emph{Against it.} Georg Christoph Lichtenberg answered in his \emph{Sudelbücher}. In notebook K, entry 76, he wrote that one should say \emph{es denkt}, ``it thinks'', as one says \emph{es blitzt}, ``it is lightning'', and that to say \emph{cogito} is already too much once it is rendered ``I think''. The indubitable datum is that thinking is occurring now. A substantial I that owns the thinking, persists and stays identical with itself is a further inference, and the datum does not supply it. Kant made the same kind of point in the Paralogisms of the \emph{Critique of Pure Reason}. The ``I'' is ``a mere consciousness that accompanies every concept'' (A346/\allowbreak{}B404). The First Paralogism (A348 to 351) is the invalid step from the formal unity of apperception to the soul as a substance; the Second and Third treat simplicity and numerical identity the same way. The ``I think'' is a logical condition on the unity of experience and reveals no substance. Both replies meet Descartes's attribute principle at the same place. The principle can be granted in general and still fail to reach this case. To apply it, one must already know that thinking, as given, is an attribute of something, and that is the point in question. The datum shows thinking and leaves unshown whether a bearer owns it. Shoemaker's and Nagel's results are results about the character and the givenness of experience, and leave the being of its subject open.

\emph{Outcome: defeat on ownership,} for the inference from indubitability to a substantial owner. The cogito delivers a self-given occurrence and no self-grounding I. The defeat has no bearing on aseity: a substantial owner, even if granted, would be a finite substance, and Descartes holds such a substance dependent (below).

\register{[PLAIN]} One more fact changes how this contest should be read. Descartes himself denies that the thinking I has aseity. In the Third Meditation he asks what preserves him in existence at each moment, finds in himself no power to conserve his own existence, and concludes that he needs a cause at each moment. So on the question this paper asks, Descartes is an \emph{ally mistaken for an enemy}: the cogito was offered as a foundation for knowledge, and its author never claimed it shows that the thinking thing exists from itself. [Argument checked 2026-10-02 at the record level (Nolan, ``The Third Meditation'', PhilArchive NOLTTM). The AT page range is UNVERIFIED.]

\subsection*{5.2 Strawson's real subject}

\register{[STEELMAN]} \emph{For a self-grounding subject.} Galen Strawson makes the strongest contemporary analytic case for a real phenomenal subject. He has argued at length against eliminativist, illusionist and deflationary accounts of experience. Two parts of his view bear directly on aseity. First, a Russellian claim: physics describes only the structure and relations of matter, and experience may be what matter is intrinsically, from the inside (Russell 1927; Strawson 2006). Second, real subjects of experience, which he calls SESMETs, subjects of experience that are single mental things (Strawson 1997, 2009). If experience is the intrinsic nature of the physical, and subjects are real things, a reader might hope to find here a subject whose being is its own.

\emph{Against it.} Strawson's view supplies both halves of the flame. Russellian intrinsicality gives the heat: character that is non-relational and real. The SESMETs give the fuel. On his own account they are short-lived; they arise, last briefly and end. By P4, a thing that begins to exist lacks aseity. The derivation is valid on the definition, and Strawson's own ontology supplies the premise. He also states the denial directly elsewhere: the first premise of his ``Basic Argument'' is that nothing can be \emph{causa sui} (Strawson 1994). The paper takes only that metaphysical premise. The argument it opens leads to an ethical conclusion, that ultimate moral responsibility is impossible, and the paper takes nothing from that conclusion (\emph{E}, kept apart).

\emph{Outcome: ally mistaken for an enemy.} Strawson never takes the step to aseity, and his SESMET view rules it out. His work establishes that experience has a real, intrinsic character, which this paper grants in full.

\register{[OPEN]} The paper cannot dispose of one part of his view. If the fundamental level of the physical world is experiential, as Strawson's panpsychist argument holds, is that fundamental level dependent? The original edition said it is ``as causally and ontologically dependent as any physical process''. That sentence asserts what the paper declines to assert for Advaita's witness. \emph{For dependence:} fundamental physical reality is described by laws and relations, and nothing in physics marks any part of it as a se. \emph{Against:} whatever is fundamental has, by definition, nothing beneath it, and an experiential fundamental level would be a candidate for an ultimate aseity. \emph{What would settle it:} an account of whether being fundamental and being dependent can coincide. The question sits at the ultimate level beside C2 and is bracketed with it. C2 asks about the witness, and this question asks about the fundamental physical level, so the paper keeps the two apart.

\subsection*{5.3 Self-luminosity and pre-reflective self-awareness}

\register{[STEELMAN]} \emph{For a self-grounding consciousness.} The Sanskrit \emph{svaprakāśa} means self-luminosity: consciousness lights itself and needs no outside light, as a lamp shows other things and needs no second lamp to be seen. Phenomenology reached a related point. Zahavi, drawing on Husserl and Sartre, argues that experience is always already given as mine, and that this givenness belongs to the structure of experiencing; no later act of reflection adds it (Zahavi 2005; Sartre 1943). If consciousness needs nothing else to be manifest, one may ask why it would need anything else to exist.

\emph{Against it.} Two Madhyamaka arguments press here. As usually cited, Nāgārjuna argues in \emph{Mūlamadhyamakakārikā} 7 that a lamp does not light itself, since there is no darkness in it to dispel. [UNVERIFIED at the verse: the attribution follows the earlier review passes of 2 October 2026; the search made for this edition did not reach the verse text.] Candrakīrti, in chapter 6 of the \emph{Madhyamakāvatāra}, turns against the Yogācāra doctrine of reflexive awareness (\emph{svasaṃvedana}). Neither targets Advaita's \emph{svaprakāśa}, which belongs to a later debate, or Zahavi. Applying their reasoning to those positions is this paper's own step. The point it carries is the one the flame already made: being manifest from itself is a fact about how a thing is known, and the lamp's burning still depends on conditions.

\emph{Outcome: ally mistaken for an enemy,} for Zahavi. His argument is phenomenological and concerns the structure of experience as lived. It makes no claim about the ultimate ontological status of the subject, so it never takes the step that the lamp argument blocks. The doctrine of self-luminosity survives in full as a claim about self-givenness.

\emph{A stronger phenomenological contrary, not examined.} Zahavi is the phenomenologist who makes no claim about ultimate status. Michel Henry does make one. On his account self-affection belongs to absolute Life, which generates itself, and each living self is generated within that Life; he identifies that Life with God (Henry, \emph{I Am the Truth}). Henry is an ally on the generated self and a contrary at the ultimate level. His ultimate claim is theological as he states it, so under BR-1 and BR-2 the paper could take up only its metaphysical content, and this edition does not.

\register{[PLAIN]} The original edition attributed the lamp argument to Candrakīrti and named his target as \emph{svaprakāśa}. The target was wrong, and this edition corrects it: Candrakīrti's target in chapter 6 is Yogācāra reflexive awareness. The lamp argument itself is Nāgārjuna's (\emph{Mūlamadhyamakakārikā} 7; UNVERIFIED at the verse). Whether Candrakīrti also uses it in \emph{Madhyamakāvatāra} 6 is unchecked. [Checked 2026-10-02: abstract of ``The concept of svasaṃvedana in Dignāga and Candrakīrti'', \emph{Asian Philosophy} 32(4), 2022, and summaries of \emph{Madhyamakāvatāra} 6.]

\subsection*{5.4 Fichte's self-positing I}

\register{[STEELMAN]} \emph{For aseity.} Fichte's \emph{Wissenschaftslehre} of 1794 to 1795 is perhaps the most systematic attempt to derive the self-grounding of the I from the structure of self-consciousness. It starts from the \emph{Tathandlung}, the act that is also a fact: the I posits itself, and this self-positing is necessary and unconditioned, the foundation of all reality. The adverb is \emph{schlechthin}: absolutely, without condition. The argument responds to two problems at once, the Kantian thing in itself and the need for a starting point that rests on no hypothesis. If the I posits itself in this sense, it has aseity. Dieter Henrich's reading strengthens this side. Henrich credits Fichte with an ``original insight'': any theory on which the self knows itself by reflecting on itself runs in a circle, since the reflecting self must already know that what it finds is itself. The \emph{Tathandlung} is Fichte's way out of that circle (Henrich 1982). A Fichtean can add a second point. The absolute I is an act, and asking whether an act ``exists from itself'' may apply a category made for substances to something that is no substance.

\emph{Against it.} Henrich's circle concerns how self-consciousness is explained. The contest here is about whether the I's existence is from itself, and on that question two points bear. First, \emph{schlechthin} declares the unconditioned status of the self-positing; the 1794 to 1795 text states it as the first principle and gives no further ground for it. Second, Fichte's own system concedes dependence at the empirical level. Determinate consciousness requires the \emph{Anstoß}, the check from the not-I, so the I anyone can identify with depends on a condition it does not give itself. The absolute I then works as a regulative ideal, a guiding idea with no concrete existence. Henrich himself argues that the 1794 formula states the insight imperfectly and that Fichte's later versions improve on it. Ernst Tugendhat pressed a related point about the presuppositions of self-acquaintance (Tugendhat 1986).

\emph{Outcome: defeat} of the claim as stated, at the strength the original edition gives it. \emph{Schlechthin} declares the step to aseity and does not earn it, and the \emph{Anstoß} concedes the empirical I's dependence. A Fichte scholar may answer that the later versions of the \emph{Wissenschaftslehre} rework exactly these points. This edition has not examined them, and the verdict covers the 1794 to 1795 text only. The verdict also leaves the act reading unanswered: if the absolute I is pure act and no existent, this paper's question may not apply to it.

\subsection*{5.5 Advaita's witness against Madhyamaka}

\register{[STEELMAN]} \emph{For aseity at the ultimate level.} The fifth case is the strongest, and it differs from the others in two ways. First, Advaita places aseity only at the ultimate level: the witness-consciousness (\emph{sākṣin}), the self-luminous Ātman identical with Brahman, the ``fourth'' (\emph{turīya}) beyond waking, dream and deep sleep. It denies aseity to the empirical ego (\emph{jīva}), which it treats as real as a phenomenon and illusory at the ultimate level (\emph{mithyā}). At the empirical level Advaita agrees with every verdict above. Second, its argument has its own shape. The witness cannot be denied without circularity, since any denial occurs within consciousness, and the consciousness that makes the denial possible cannot be made its object without regress. The witness is the ground of all illumination, so it cannot be derived from anything without the relation of derivation going unexplained. A third argument targets unconditionedness directly. Across waking, dream and deep sleep the contents of experience come and go, and the witness is present and unchanged through all three, so it is conditioned by none of them; the three states point to the fourth, \emph{turīya} (IEP, ``Advaita Vedānta''; the argument as a named locus in Gauḍapāda or Śaṅkara is UNVERIFIED). From ineliminability, non-objectifiability and invariance through the three states, Śaṅkara, in the tradition that follows him, concludes that the witness is self-established. [The original edition said ``the tradition he founded''; a reviewer noted, from memory, that Gauḍapāda's \emph{Kārikā} is earlier (UNVERIFIED).]

The paper reconstructs that last step as a grounding excluded middle, P6: what cannot be derived and cannot be denied must be self-grounded, because everything is either other-grounded or self-grounded. This is the reconstruction made in the adversarial analysis behind the paper; no passage of Śaṅkara is cited for it. Gupta (1998), the standard study of the \emph{sākṣin}, is the place to look for one. An Advaitin may reject this idiom. The witness is \emph{svataḥ-siddha}, self-established, and Brahman is often said to stand outside the relation of ground and grounded altogether, so ``self-grounded'' may misdescribe the claim. The reconstruction is the paper's, and the excluded middle below is the paper's statement of it. The original edition calls the ineliminability argument ``sound''. This edition grants its premise, that any attempt to eliminate or derive the witness presupposes it, and does not argue for it.

\emph{Against it.} Madhyamaka denies the excluded middle. Things arise in dependence on conditions (\emph{pratītyasamutpāda}), and no ground, self or other, stands beneath them free of dependence. That is a third option: groundless dependent arising, emptiness (\emph{śūnyatā}). If it is coherent, as contemporary scholarship on Nāgārjuna argues (for example, Westerhoff 2009), then ineliminability and non-derivability do not yield self-establishment. The witness could be ineliminable from every act of awareness and still be empty of existence from itself. The same holds for invariance: a witness unchanged through waking, dream and sleep is shown to be unconditioned by those states, and that falls short of being conditioned by nothing. On this reading the appearance of aseity is a deep case of grasping at intrinsic existence (\emph{svabhāva-graha}).

\emph{Advaita's reply.} If nothing grounds dependent arising, what sustains it? A chain of dependently arisen things with no ground risks an infinite regress or a brute, unexplained fact of arising. Advaita argues that such a chain cannot be coherent without a self-luminous witness as its condition.

\emph{Madhyamaka's rejoinder.} The question ``what sustains it?'' assumes that dependent arising needs a sustainer, and that assumes the things arising have an intrinsic nature (\emph{svabhāva}) that must be held up by something. Emptiness denies exactly that. Many Mādhyamikas would add that they advance no thesis of their own to be grounded (\emph{Vigrahavyāvartanī} 29; UNVERIFIED at the verse). On this view the regress worry is a symptom of the grasping the analysis removes.

\emph{Outcome: standoff.} Each side has a strong argument and an unpaid debt. Madhyamaka must show that groundlessness is coherent when applied to the witnessing awareness itself. Advaita must show that its excluded middle is more than an assumption. Neither has defeated the other in the literature, and this paper does not try to decide between them. The standoff sits entirely at the ultimate level. Advaita agrees that the empirical ego has no aseity.

\subsection*{5.6 Albahari's ownerless consciousness: a contrary added in this edition}

\register{[STEELMAN]} \emph{The contrary.} Miri Albahari argues that the self, as a bounded owner of experiences and a controlling agent, is an illusion, and that the illusion is built from two tiers: a tier of non-illusory awareness, and a tier of desire-driven thought and emotion (Albahari 2006). Many properties mistakenly credited to the self, she argues, belong to that awareness: it is unified, unbroken, invariable and unconstructed. In nirvana, on her reading, awareness would be pure, without an intentional object, and unconditioned by space and time (Albahari 2011). [Positions checked 2026-10-02 from the publisher's description of the 2006 book and the PhilPapers record of the 2011 chapter; neither text was read.]

\emph{Why it matters here.} The original edition defined seity as an ownerless occurrence, self-given and dependent, and treated ``ownerless'' and ``dependent'' as going together: dissolve the owner and dependence remains. Albahari's position separates them. Her conclusion about the owner matches Lichtenberg's, since no owner is given, and she still holds that the awareness left over is unconditioned. Her view is ownerless and, at its limit, unconditioned. Whether her ``unconditioned'' (Pali \emph{asaṅkhata}) amounts to aseity in this paper's sense is the paper's question, and a reader of Albahari may decline to call it her position.

\emph{Outcome: ally on one question; on the other, reported and not contested.} On ownership she is an ally: the owner is constructed, which is the paper's reading of \emph{es denkt}. On aseity this edition states her conclusions from records and gives neither her argument nor a reply, so no contest has taken place. Her view sits at the ultimate level beside Advaita's witness, where the paper's bracket already stands. Her view also corrects the paper. Dissolving the owner does not deliver dependence, so this edition defines seity as self-given, occurring and not a se, and treats ownerlessness as a separate result that Lichtenberg and Kant support for the empirical datum. Under that definition Strawson's case still counts: his subjects are owners of a thin kind, and they are generated, so they lack aseity.

\section*{6. The premise that carries dependence}

\register{[STEELMAN]} The contests in section 5 remove arguments \emph{for} aseity. Failing to establish aseity is weaker than establishing dependence. The original edition took the second step without saying which premise carried it. This section sets that premise against its strongest contrary.

\emph{For dependence, the generation route.} The everyday first-person subject, taken as the embodied person of one life, began (P4a). Before a certain time there was no such person, and its experiences have a developmental history. By P4, a thing that begins to exist lacks aseity. The derivation is valid. Not every party grants its premise for every sense of ``subject''. Advaita holds that the existence of the individuated \emph{jīva} is without a beginning (IEP, ``Advaita Vedānta''), and Buddhist traditions hold the mental continuum beginningless (UNVERIFIED here; Nāgārjuna reports the Buddha's teaching that no prior limit of saṃsāra is known, that it has no beginning or end, \emph{Mūlamadhyamakakārikā} 11.1). Advaita still denies the \emph{jīva} aseity, on other grounds: it is \emph{mithyā}, dependent on ignorance. So Advaita reaches the paper's conclusion for the individuated self by a different premise. It is an ally on the conclusion only.

\emph{For dependence, the conditioning route.} The occurrence of that subject's experience is causally conditioned by neurological processes, development and ongoing causal relations (P5). Experience varies with the state of the brain, and its occurrence tracks conditions the subject does not provide. [UNVERIFIED as stated: the original edition calls this ``the ordinary finding of the sciences of mind'' with no citation, and none was opened for this edition.] To reach a metaphysical conclusion, this route needs a bridge:

\begin{itemize}
\item \textbf{BR-3.} If a thing's occurrence is causally conditioned by other things, its existence is not from itself. This holds at the empirical level only, and makes no claim about whatever conditions the empirical order as a whole. A reader may reject BR-3 on its own.
\end{itemize}
\emph{Against it.} A dualist, an idealist and an Advaitin can all accept that experience correlates with neural states and varies with them. Each denies that this settles existential dependence. For the dualist, the brain conditions what a mind does; for the idealist, the brain is itself an appearance within experience; for the Advaitin, conditioning touches the ego and leaves the witness untouched. Against the generation route, a transcendental philosopher can reply that the subject in question is no item in time at all, so ``began'' does not apply to it. In each reply, what the evidence shows to be conditioned or generated is the empirical subject, and the subject that perspectival aseity is claimed for lies elsewhere.

\emph{Outcome: survival by concession.} For the embodied person of one life, the generation route stands: that person began, and P4 then gives the result by definition. The conditioning route is exactly as strong as BR-3. Every reply above works by moving the subject of the claim out of the empirical order, into a mind distinct from the brain, a transcendental subject, or the witness. That move concedes the empirical case and carries the dispute to the level the bracket already covers.

\register{[OPEN]} Does causal conditioning of occurrence amount to dependence for existence?

\emph{Yes:} a thing that exists only while certain conditions hold, and ceases when they fail, has its existence from those conditions in every sense that matters to aseity; an a se being has nothing whose failure could end it. \emph{No:} conditioning may govern what appears or how it appears without governing the being of what appears; a radio's output depends on the circuit, and the broadcast does not. This is the transmission or filter view William James set out, on which the brain admits and shapes consciousness and does not produce it (James 1898). \emph{What would settle it:} an account of existential dependence precise enough to say whether the sustaining conditions of an occurrence are conditions of its existence, together with an argument that the subject of experience is the occurrence itself and is no further thing that the occurrence shows. The SEP entry by Tahko and Lowe sets out the candidate relations. The paper does not choose among them, and C1's dependence form rests on the answer.

\register{[PLAIN]} Part of this question can be sharpened. On the SEP definition, a thing depends rigidly for its existence on another when ``necessarily, x exists only if y exists'', and the necessity meant is metaphysical (Tahko and Lowe). The evidence about neural conditioning supplies, at most, necessity under the laws of nature. So the open question reduces to a narrower one: does necessitation under natural law suffice for existential dependence?

\section*{7. What the empirical result is, stated plainly}

\register{[PLAIN]} Three things can be said without hedging.

\begin{enumerate}
\item No argument examined in this paper establishes aseity for the everyday first-person subject. In the cogito, the step from indubitability to a substantial owner fails. In Strawson and Zahavi, the step to aseity is never taken. In Fichte, \emph{schlechthin} declares it for the absolute I, and the \emph{Anstoß} concedes the empirical I's dependence. Advaita does not claim it at the empirical level. This is a report on the arguments examined, and it holds whatever the reader decides in section 6.
\item The embodied person of one life, which began, lacks aseity by the definition of aseity (P4a, P4). This is a valid derivation. Its premise holds for the person of one life; Advaita holds the individuated self beginningless, rejects the premise for that self, and denies that self aseity on other grounds (section 6).
\item The original edition's label for its empirical conclusion, ``decisive (high confidence)'', was stronger than the supports it cited. The supporting diagnosis was advanced at moderate-to-high confidence, and the premise carrying dependence was asserted without argument. This edition withdraws that label from the dependence claim and states the result in the two parts above.
\end{enumerate}
What does not follow plainly is that the subject of experience, in whatever sense a defender of perspectival aseity means, is existentially dependent. That is the open question of section 6.

\section*{8. The convergence of the critiques}

\register{[STEELMAN]} \emph{For the convergence.} The original edition argued that three critiques from different traditions reach the same diagnosis: Lichtenberg's \emph{es denkt}, Kant's Paralogisms, and the Madhyamaka critique of self-illumination. They share no authority and no method. When such traditions agree, their agreement weighs more than any one of them.

\emph{Against it.} The three critiques target different claims. Lichtenberg targets ownership. Kant targets the inference to a soul as substance. The Madhyamaka arguments target reflexive awareness. None of them argues that the existence of the subject depends on conditions. The attribution of the third critique was also wrong in the original edition (5.3), and its application to Zahavi and Advaita is the paper's own step. And all three sit inside one problem, the analysis of the first-person standpoint. Agreement within one domain is weaker evidence than agreement reached from a different domain entirely, as the original edition itself said.

\emph{Outcome: survival by concession.} The convergence survives in a narrower form: none of the three finds a self-grounding owner in self-givenness. That is significant supporting evidence for the result on ownership, at moderate-to-high confidence, the label the original edition gave it in its own analysis. It is no evidence for dependence.

\section*{9. The no-aseity thesis at its terminus}

\register{[PLAIN]} The universal no-aseity thesis, ``no aseity anywhere, not even at the ultimate floor'', cannot be maintained on this paper's evidence. Here ``anywhere'' covers the ultimate ground of experiencing, and the verdict bears on no theological thesis (BR-1). Advaita's witness stands undefeated at the ultimate level, and Albahari's unconditioned awareness, reported in 5.6 and not examined there, is a further candidate. An undefeated contrary leaves the universal thesis unestablished and leaves its falsity unshown. The paper therefore brackets the universal form, as the original edition did, and whether it is true is the open question below. The bracketed form survives: within the conferred order, no argument examined establishes aseity for the first-person subject.

\register{[OPEN]} Is the ultimate ground of experiencing a self-luminous aseity, or is even that ground empty of existence from itself?

\emph{For self-luminous aseity:} the witness is presupposed by every act of awareness, including any act that would deny it or derive it; it stays unchanged through waking, dream and deep sleep while their contents pass; groundless dependent arising risks regress or brute fact (Advaita); pure awareness has properties, unity, invariability, unconstructedness, that no conditioned process seems to share (Albahari). \emph{For emptiness:} ineliminability from awareness is a fact about how awareness is structured, and gives no reason to think the witness exists from itself; the excluded middle that converts it into self-establishment is assumed without being shown; a groundless chain is no worse off than a ground that must explain itself; the demand for a sustainer assumes the intrinsic nature that emptiness denies, and Madhyamaka offers no rival thesis to be grounded (Madhyamaka). \emph{What would settle it:} an argument that the grounding excluded middle holds, which would favour Advaita; or a coherent account of groundless dependence applied to awareness itself, which would favour Madhyamaka. Neither exists in a form both sides accept. The paper brackets the question because neither side has defeated the other, and it leaves the question with you.

\section*{10. Further objections}

\subsection*{10.1 Who holds the view?}

\register{[PLAIN]} No living philosopher was found who defends aseity for the everyday first-person subject. The search behind that null was limited: the literature records opened for this edition on 2026-10-02, around the five cases and Albahari, and no systematic survey. Of the five original cases, Descartes, Strawson and Zahavi decline the step, Fichte makes it for an absolute I, and Advaita makes it for the ultimate witness only. Candidates a reader may raise sit at the ultimate level only: Michel Henry (5.3) and cosmopsychism (Goff 2017; UNVERIFIED). So the paper is best read as a test of the arguments that would support empirical perspectival aseity.

\register{[OPEN]} Is empirical perspectival aseity a live position at all? \emph{Yes:} the pull it describes is real, and many readers feel it before any argument; a test of the arguments for it serves those readers. \emph{No:} at the empirical level the view is a temptation more than a held position, and if no one holds it, the paper's empirical result is a clarification of what its sources already say. \emph{What would settle it:} a named defender, or a reason to think the view operates without one, for instance as an assumption inside a wider position.

\subsection*{10.2 Does this support eliminativism or illusionism?}

\register{[STEELMAN]} \emph{Objection.} An argument that removes the self-grounding subject helps eliminativism, illusionism or a naive no-self view.

\emph{Reply.} The paper grants the self-givenness and intrinsic character of experience throughout and disputes only the ontological conclusion drawn from them. What remains after the dissolution is a real occurrence, self-given and not a se. So the paper gives eliminativism no help; it accepts the phenomena eliminativism denies.

\emph{The illusionist's objection.} Keith Frankish's illusionism holds that phenomenal consciousness, as ordinarily conceived, is an introspective illusion (Frankish 2016). An illusionist would press a different point: the for-me-ness the paper grants in section 2 is the claim in dispute, and the paper assumes it without argument. Guillot adds that arguments from subjective character often slide between for-me-ness, me-ness and mineness (Guillot 2017).

\emph{Reply.} The paper's thesis is conditional on that premise. It argues that even if self-givenness is granted at full strength, aseity does not follow. If the illusionist is right, the argument for perspectival aseity loses its starting point, and the paper's conclusion holds a fortiori. The paper does not defend the premise against illusionism, and section 2 says so.

\emph{Outcome: ally mistaken for an enemy, on the conclusion only.} If illusionism is true, no self-grounding subject is given, which is the paper's conclusion on ownership. Frankish does not write about aseity, so the shared conclusion is this paper's inference from his view. On the premise, whether experience is as it seems, the two disagree, and the paper does not settle that disagreement.

\subsection*{10.3 The essential indexical}

\register{[STEELMAN]} \emph{Objection.} John Perry's essential indexical shows that the first person is more than a dependent occurrence (Perry 1979).

\emph{Reply.} Perry shows that first-person reference cannot be replaced by third-person description. That irreducibility belongs to how the perspective is given. It is a self-givenness result, and it is compatible with the dependent existence of the perspective so given.

\emph{Outcome: defeat} of the objection, as an argument for aseity. The indexical result stands untouched as a result about givenness.

\section*{11. Update, October 2026}

\emph{This section was added on 2026-10-02 at the author's direction. It stands apart from the argument above and changes no step in it. Its claims are proposed, and each carries a domain tag.}

\textbf{\register{[PLAIN]} What Jaynes claimed.} In \emph{The Origin of Consciousness in the Breakdown of the Bicameral Mind} (Houghton Mifflin, 1976), Julian Jaynes argued that consciousness, in his narrow sense of introspection and a narrating ``I'', is learned through language and culture and is not innate. He placed its arrival late, in the second millennium BCE. Before it, on his account, people heard their own cognition as commanding voices they took to be gods. \emph{C-E.} [Checked 2026-10-02 at the level of encyclopedia and bibliographic records; the 1976 text was not opened, and no page is cited until it is.]

\textbf{\register{[STEELMAN]} Jaynes against his critics.} Ned Block's review argued that Jaynes confuses the nature of people's thought processes with the nature of their theories about them, and pointed to planning and deceit in chimpanzees and in the ancient texts Jaynes reads (Block 1981). Defenders reply that Block substitutes a broader sense of ``consciousness'', close to perception, for Jaynes's narrow sense of introspection; Gary Williams makes this defence at length (Williams 2011). Daniel Dennett gave the book its most sympathetic philosophical reading, treating Jaynes's thesis as a claim about the software of the mind, a learned way of organising cognition, as distinct from its hardware (Dennett 1986; UNVERIFIED this session). Cavanna and colleagues reappraised the hypothesis thirty years on, naming its neurological basis and the philological accuracy of its readings as the main critical issues (Cavanna et al. 2007). The dating and the right-hemisphere hypothesis remain the most contested parts of the book. \emph{Outcome: standoff on the narrow thesis.} Whether introspective selfhood is learned is an empirical question this paper cannot settle. Block's point stands as a warning against reading Jaynes as a claim about experience itself.

\textbf{\register{[PLAIN]} Different targets.} Jaynes's consciousness is reflective, narrated selfhood. Self-givenness in this paper is pre-reflective and phenomenal. The two must be kept apart, and Block's critique is the reason.

\textbf{The author's points.} The author draws three points, paraphrased here from his notes of 2 October 2026. Whatever a person holds to be the source of their life or consciousness, God, a mother or another being, no one can argue they came from nothing (\emph{C-M, about persons}). Even prophecy, a word received from God, formed its new discourse only when a second stream of thought, divine or human, joined it (\emph{T, with a historical claim inside it}). Some things can be found alone, and some only in a meeting of minds (\emph{X, a claim about inquiry}).

\textbf{\register{[PLAIN]} How far these points reach.} The first point can be read two ways, and the author has not yet chosen between them. Read as stated, ``no one came from nothing'' asserts that every person has a source. That is a claim of dependence, and it sides with the ``Yes'' answer to the open question of section 6. Read more narrowly, it restates the paper's P4 for persons: whatever one's account of a person's source, the person began. The case that tests the point across belief systems is the Advaitin's Ātman, which on that view never came to be at all. The theological points are no evidence for either reading. They show only that the point does not depend on any one view of where a person comes from. Jaynes's thesis and the theological points are no evidence for each other: heard voices in Jaynes's account are a claim about cognition, and say nothing for or against prophecy. No ethical conclusion is drawn from any of them.

\begin{itemize}
\item \textbf{Bridge BR-J.} If the owned, narrating ``I'' is acquired through language and culture, then the owner the cogito names is a received construction. That reaches a conclusion close to Lichtenberg's \emph{es denkt} by a different route: Lichtenberg's point concerns what the datum licenses, Jaynes's the history of a capacity. \emph{C-E to C-M, through BR-J, which a reader may reject on its own.} This bears on ownership only. It says nothing about whether experience depends on conditions, and nothing about the ultimate level.
\end{itemize}
\register{[OPEN]} If the narrating ``I'' was learned, what was there before it learned to say ``I''? \emph{One reading:} a self-given stream not owned by an ``I'', and heard as another's, which is close to seity in this paper's sense, and evidence that the owner was always a construction. \emph{Another:} Jaynes describes the loss and gain of a concept, and experience, owned or not, went on unchanged beneath it, so his history bears on the concept of a self and leaves the self-givenness of experience where it was. \emph{What would settle it:} a way to tell, from the historical and clinical record, a change in how experience is described from a change in how it is had. The author's third point, that some discoveries need a meeting of minds, has no stated bearing on this paper.

\section*{12. Questions for the reader}

\register{[OPEN]} The paper ends with the questions it could not close. Each is set out above with the considerations on both sides.

\begin{enumerate}
\item Is the ultimate ground of experiencing a self-luminous aseity, or empty of existence from itself? (Section 9)
\item Does causal conditioning of an occurrence amount to dependence for existence? (Section 6)
\item Can an ownerless awareness be unconditioned, as Albahari holds? If it can, what is left of the intuition that a subject without an owner must be a passing event? (Section 5.6)
\item Is the fundamental level of a panpsychist world dependent? (Section 5.2)
\item Does the fire and fuel image survive Nāgārjuna's chapter 10 as anything more than a conventional picture? (Section 3.3)
\item Is empirical perspectival aseity a position anyone holds, or a temptation that precedes argument? (Section 10.1)
\item What did the narrating ``I'' replace, if Jaynes is right that it was learned? (Section 11)
\end{enumerate}
You may reach answers the author has not. The paper's own verdicts are labelled so that you can disagree with each one separately.

\section*{13. What can be said plainly}

\register{[PLAIN]}

\begin{itemize}
\item Self-givenness and self-grounding are different properties, and intrinsic character does not entail aseity.
\item No argument examined here establishes aseity for the everyday first-person subject.
\item The embodied person of one life began, and so lacks aseity by the definition of aseity. Advaita, which holds the individuated self beginningless, denies it aseity on other grounds.
\item The cogito's inference of a substantial owner from indubitability fails, as Lichtenberg and Kant showed. That result concerns ownership and leaves aseity where Descartes left it, denied.
\item Descartes, Strawson and Zahavi decline the step to aseity themselves.
\item The universal no-aseity thesis cannot be maintained on this paper's evidence, and the paper brackets it. Advaita's witness stands undefeated at the ultimate level.
\item The original edition named the wrong target for the lamp argument, overstated the label on its empirical conclusion, and bundled ownerlessness with dependence. This edition corrects all three and says so.
\item Nothing here bears on God, on religious truth, or on what anyone ought to do (BR-1, BR-2).
\end{itemize}
\section*{Reader's note on AI assistance}

This paper was drafted with AI assistance. The question it answers, and the willingness to let the thesis fail, came from the author in June 2026. The five traditions, the distinction between self-givenness and self-grounding, the flame test, the reconstruction of each tradition and every verdict came from AI model runs in an adversarial analysis that the author directed. This edition, also AI-assisted and written on 2026-10-02, restores the original edition's structure of contests between the thesis and its contraries, adds Albahari's view as a sixth contrary, and changes several conclusions. The draft was revised the same day under an independent review. The author has not yet re-derived each claim from the primary sources or checked every citation himself. Until he does, the paper is research material and cannot yet count as original scholarship. The notice below is kept as written.

\begin{quote}
\textbf{Substrate Notice.} This manuscript is AI-assisted research substrate produced from a frozen adversarial-analysis file (ASEITY-CONFRONTATION.md) and related corpus materials. It is not submittable as original human work until the candidate has re-derived and takes intellectual ownership of each claim from primary sources, verified all citations independently, and satisfied applicable institutional AI-use policies. Re-derivation means what dissertation/\allowbreak{}00-CONTRIBUTION.md and the curation require: working each argument cold from the primary texts (Descartes, Lichtenberg, Kant’s First Critique, Candrakīrti, Śaṅkara, the Madhyamaka literature) and not paraphrasing this draft. The philosophical positions stated here are bounded by the adversarial analysis on which they are based. The Sources below are not yet reconciled to the corpus’s canonical bibliography (references/\allowbreak{}BIBLIOGRAPHY.md), which permits citing only from its approved list and prohibits asserting page numbers that have not been independently verified. Every venue, volume, and page specific below is therefore provisional and must be verified against the edition used before placement.
\end{quote}

\section*{References}

Status codes record what was done on 2026-10-02 for this edition. \textbf{Opened} means a bibliographic or summary record was read this session at the source named; it does not mean the work itself was read. \textbf{UNVERIFIED} means no record was opened this session (some were checked in an earlier pass on 2 October 2026, as noted). Ranges use ``to'' in place of dashes.

\begin{itemize}
\item Albahari, Miri. \emph{Analytical Buddhism: The Two-Tiered Illusion of Self}. Palgrave Macmillan, 2006. \emph{Opened: publisher and repository records.}
\item Albahari, Miri. ``Nirvana and Ownerless Consciousness.'' In Siderits, Thompson and Zahavi (eds.), \emph{Self, No Self?}, Oxford UP, 2011. \emph{Opened: PhilPapers record ALBNAO and summary.}
\item Aquinas, Thomas. \emph{Summa Theologiae} I, q. 2, a. 3. \emph{Opened: CCEL text via search.} The article denies that a thing can be the efficient cause of itself; the standard English reads ``neither is it, indeed, possible'', the Latin \emph{nec est possibile}. Not cited in the body; consulted for the correction that drops the original edition's sentence saying the scholastics called self-grounding \emph{causa sui}.
\item Block, Ned. Review of Jaynes. \emph{Cognition and Brain Theory} 4 (1981): 81 to 83. \emph{Opened: PhilPapers record BLOROJ-2 and summaries; the review itself was not read.} The summary in section 11 was corrected after a later reference audit read pages 81 to 83.
\item Candrakīrti. \emph{Madhyamakāvatāra}, chapter 6. \emph{Opened: abstract of ``The concept of svasaṃvedana in Dignāga and Candrakīrti'', Asian Philosophy 32(4), 2022, and text summaries.} Editions (La Vallée Poussin; Huntington 1989) UNVERIFIED.
\item Cavanna, A. E., M. Trimble, F. Cinti and F. Monaco. ``The `Bicameral Mind' 30 Years On.'' \emph{Functional Neurology} 22(1) (2007): 11 to 15. \emph{Opened: PubMed listing 17509238 (title only).} The summary in section 11 was corrected after a later reference audit read the abstract.
\item Dennett, Daniel C. ``Julian Jaynes's Software Archeology.'' \emph{Canadian Psychology} 27(2) (1986): 149 to 154. \emph{UNVERIFIED this session (checked in an earlier pass on 2 October 2026).}
\item Descartes, René. \emph{Meditationes de Prima Philosophia} (1641), Second and Third Meditations. \emph{Opened for the Third Meditation preservation argument: Nolan, PhilArchive NOLTTM; AT pages UNVERIFIED.}
\item Descartes, René. \emph{Principia Philosophiae} (1644), I.11 and I.52. \emph{UNVERIFIED.}
\item Fichte, J. G. \emph{Grundlage der gesamten Wissenschaftslehre} (1794 to 1795); trans. Heath and Lachs. \emph{UNVERIFIED this session (checked in an earlier pass on 2 October 2026).}
\item Frankish, Keith. ``Illusionism as a Theory of Consciousness.'' \emph{Journal of Consciousness Studies} 23(11 to 12) (2016): 11 to 39. \emph{Opened: PhilPapers FRAIAA-4.}
\item Goff, Philip. \emph{Consciousness and Fundamental Reality}. Oxford UP, 2017. \emph{UNVERIFIED.}
\item Guillot, Marie. ``I Me Mine: On a Confusion Concerning the Subjective Character of Experience.'' \emph{Review of Philosophy and Psychology} 8(1) (2017): 23 to 53. \emph{Opened: search records (PhilPapers GUIIMM, PubMed 28386303, Springer DOI 10.1007/\allowbreak{}s13164-016-0313-4).}
\item Gupta, Bina. \emph{The Disinterested Witness: A Fragment of Advaita Vedānta Phenomenology}. Northwestern UP, 1998. \emph{Opened: PhilPapers and Google Books records.}
\item Henrich, Dieter. ``Fichte's Original Insight,'' trans. D. Lachterman. \emph{Contemporary German Philosophy} 1 (1982): 15 to 53. \emph{Opened: search records (Wikipedia ``Dieter Henrich''; course PDF listing at phil880.colinmclear.net; Frank, PhilPapers FRAFFO) for the reflection-circle reading credited to Fichte.} German original 1966 UNVERIFIED.
\item Henry, Michel. \emph{I Am the Truth: Toward a Philosophy of Christianity}, trans. Susan Emanuel. Stanford UP, 2003. \emph{Opened: search records (SEP ``Michel Henry'' listing, PhilArchive HAOIMH, publisher listings) for Life's self-generation and its identity with God; text not read; translation year UNVERIFIED.}
\item IEP, ``Advaita Vedānta'' (iep.utm.edu/\allowbreak{}advaita-vedanta). \emph{Opened: ``The existence of an individuated jīva and the world are without a beginning''; the three states ``point to'' turīya.}
\item James, William. \emph{Human Immortality: Two Supposed Objections to the Doctrine}, 1898. \emph{Opened: search records (University at Buffalo course PDF listing and summaries) for the transmissive function of the brain.}
\item Jaynes, Julian. \emph{The Origin of Consciousness in the Breakdown of the Bicameral Mind}. Houghton Mifflin, 1976. \emph{Opened: bibliographic and encyclopedia records; text not read.}
\item Kant, Immanuel. \emph{Kritik der reinen Vernunft}, A 1781, B 1787. A346/\allowbreak{}B404 for ``a mere consciousness that accompanies every concept''; First Paralogism A348 to 351. \emph{Opened: Forgione, PhilArchive FORKAT-6 (quotation); IEP ``Kant: Philosophy of Mind'' (numbering). The quoted wording is that of the Guyer and Wood translation (Cambridge University Press, 1998); translation in hand UNVERIFIED.}
\item Lichtenberg, G. C. \emph{Sudelbücher}, notebook K, entry 76. \emph{Opened: zeno.org and quotation records of K 76.} Hollingdale translation UNVERIFIED.
\item MacKenzie, Matthew D. ``The Illumination of Consciousness.'' \emph{Philosophy East and West} 57(1) (2007): 40 to 62. \emph{Opened: JSTOR 4488075 record.} Proposed reading; not cited in the body.
\item Nagel, Thomas. ``What Is It Like to Be a Bat?'' \emph{Philosophical Review} 83(4) (1974): 435 to 450. \emph{UNVERIFIED this session (checked in an earlier pass).}
\item Nāgārjuna. \emph{Mūlamadhyamakakārikā}, chapters 7, 10 and 11. \emph{Opened for chapter 10: chapter title and summaries. Chapter 7 lamp verse UNVERIFIED.} Verse 11.1 is cited from a later reference audit that read the Sanskrit text. Edition to be chosen.
\item Nāgārjuna. \emph{Vigrahavyāvartanī} 29. \emph{UNVERIFIED.}
\item Perry, John. ``The Problem of the Essential Indexical.'' \emph{Noûs} 13(1) (1979): 3 to 21. \emph{UNVERIFIED this session (checked in an earlier pass on 2 October 2026).}
\item Russell, Bertrand. \emph{The Analysis of Matter}, 1927. \emph{UNVERIFIED.}
\item Sartre, Jean-Paul. \emph{L'Être et le Néant}, 1943. \emph{UNVERIFIED.}
\item Schaffer, Jonathan. ``On What Grounds What.'' In Chalmers, Manley and Wasserman (eds.), \emph{Metametaphysics}, Oxford UP, 2009, 347 to 383. \emph{Opened: search records; grounding characterised as irreflexive.}
\item Shoemaker, Sydney. ``Self-Reference and Self-Awareness.'' \emph{Journal of Philosophy} 65(19) (1968): 555 to 567. \emph{UNVERIFIED this session (checked in an earlier pass on 2 October 2026).}
\item Siderits, Mark, Evan Thompson and Dan Zahavi (eds.). \emph{Self, No Self?} Oxford UP, 2011. \emph{Opened: PhilPapers SIDSNS.}
\item Strawson, Galen. ``The Self'' (1997), \emph{Journal of Consciousness Studies} 4(5 to 6): 405 to 428; ``Realistic Monism'' (2006); \emph{Selves} (Oxford UP, 2009). \emph{Opened for the SESMET expansion: PhilPapers STRTSA-22 and Strawson's text. Volume and page details UNVERIFIED this session.}
\item Strawson, Galen. ``The Impossibility of Moral Responsibility.'' \emph{Philosophical Studies} 75 (1994): 5 to 24. \emph{Opened: search records of the Basic Argument (premise: nothing can be causa sui); only the metaphysical premise is used.}
\item Tahko, Tuomas E., and E. J. Lowe. ``Ontological Dependence.'' \emph{Stanford Encyclopedia of Philosophy} (substantive revision 1 May 2025). \emph{Opened: the SEP entry, section 1 (definition of rigid existential dependence; metaphysical modality).}
\item Tugendhat, Ernst. \emph{Self-Consciousness and Self-Determination}, trans. Paul Stern. MIT Press, 1986. \emph{UNVERIFIED this session (checked in an earlier pass on 2 October 2026).}
\item Westerhoff, Jan. \emph{Nāgārjuna's Madhyamaka: A Philosophical Introduction}. Oxford UP, 2009. \emph{Opened: PhilPapers WESNMA-2 and H-Net review record.}
\item Williams, Gary. ``What Is It Like to Be Nonconscious? A Defense of Julian Jaynes.'' \emph{Phenomenology and the Cognitive Sciences} 10(2) (2011): 217 to 239. \emph{Opened: PhilArchive WILWII-3 record and the Julian Jaynes Society page that cites it.}
\item Zahavi, Dan. \emph{Subjectivity and Selfhood}. MIT Press, 2005. \emph{Opened: PhilPapers ZAHSAS.}
\end{itemize}
\vfill{\footnotesize Typeset with LaTeX from the approved source \texttt{writing/papers/self-given.md}. The web page is rendered from the same source.\newline SHA-256 of the source: \texttt{50cc0e469edb09a4538ed31f7f06ee6e30d57fa9dfef08e4b94a385f862dff39}\par}
\end{document}
